\chapter{Implementation}
\label{ch:implementation}

This chapter describes the code that implements Chapter~\ref{ch:methodology}. It is organised as four work packages (data, environment, experiments, reporting). Nothing here claims engineering novelty; the scientific claim is Equation~\ref{eq:probabilistic_trade_size} --- how the confidence signal changes trade size.

\section{Work-package summary}
\label{sec:work_package_summary}

Table~\ref{tab:work_packages} maps the four work packages onto the source files in the public repository. Each task is something a reader can verify in the repository.

\begin{table}[h]
\centering
\caption{Work-package structure of the implementation. Each task points to a specific file or notebook in the public repository.}
\label{tab:work_packages}
\footnotesize
\begin{tabularx}{\textwidth}{p{2.6cm}p{4.0cm}X}
\toprule
\textbf{Work package} & \textbf{What it does} & \textbf{Deliverable artefact} \\
\midrule
\textbf{WP1 --- Data and forecaster} & Fetch daily adjusted-close prices for the market sample. Train the probabilistic LSTM in aleatoric and epistemic modes. & \path{experiments/common.py} (data + forecaster utilities); \path{experiments/runners/run_probabilistic_agent.py} (forecaster training and uncertainty estimation) \\
\midrule
\textbf{WP2 --- Environment and policy} & Implement the Gymnasium-compatible trading environment. Train PPO under three uncertainty regimes (none, aleatoric, epistemic). & \path{experiments/common.py} (\texttt{StockEnv} class); \path{experiments/runners/run_baseline.py} (Arm A); \path{experiments/runners/run_probabilistic_agent.py} (Arms B and C) \\
\midrule
\textbf{WP3 --- Experimental grid} & Run the four-agent comparison on SPY and on the 70-stock market sample. Run the four-fold walk-forward grid. & \path{experiments/runners/run_walk_forward.py}; \path{experiments/runners/run_extended_grid.py}; \path{notebooks/02_Full_Experiments.ipynb} \\
\midrule
\textbf{WP4 --- Evaluation and reporting} & Aggregate per-cell JSON results into median + IQR summaries. Render headline tables and equity curves. & \path{experiments/aggregate_results.py}; \path{reports/builders/plot_phase2_charts.py}; \path{notebooks/01_Project_Walkthrough.ipynb} \\
\bottomrule
\end{tabularx}
\end{table}

\section{Software stack}
\label{sec:software_stack}

The implementation is Python 3.11. \textbf{PyTorch} runs the forecaster neural network. \textbf{Stable-Baselines3}~\parencite{raffin2021sb3} is an off-the-shelf PPO implementation --- the learning algorithm is not reimplemented here. \textbf{Gymnasium} is the standard interface between agent and simulated market. Prices come from \textbf{yfinance}; tables and plots use \textbf{pandas} and \textbf{matplotlib}. Versions are pinned in \texttt{requirements.txt}.

Figure~\ref{fig:pipeline} shows how these pieces fit together end to end, from raw prices through the LSTM forecaster and the uncertainty injection into \texttt{StockEnv}, to PPO training and evaluation. The baseline arm skips the forecaster and uncertainty-injection stages and is otherwise identical.

% Figure: End-to-end training pipeline
\begin{figure}[htbp]
\centering
\begin{tikzpicture}[
  box/.style={draw, rounded corners=3pt, minimum height=1.2cm,
              text width=#1, align=center, font=\small},
  arr/.style={-{Stealth[length=3mm]}, thick},
]
  % Row 1
  \node[box=3cm, fill=gray!10] (prices)
    {\textcircled{1} Prices\\Yahoo daily closes};
  \node[box=3.5cm, fill=blue!10, right=0.5cm of prices] (lstm)
    {\textcircled{2} LSTM forecaster\\Gaussian NLL loss\\$\to$ uncertainty $u_t$};
  \node[box=4cm, fill=purple!10, right=0.5cm of lstm] (inject)
    {\textcircled{3} Inject $u_t$ into\\StockEnv\\(baseline: $u_t = 0$)};

  \draw[arr] (prices.east) -- (lstm.west);
  \draw[arr] (lstm.east)   -- (inject.west);

  % Row 2
  \node[box=10cm, fill=yellow!8, below=1.2cm of lstm] (ppo)
    {\textcircled{4} Train PPO (same settings both arms)\\
     lr $3 \times 10^{-4}$ $\cdot$ n\_steps 512 $\cdot$ batch 64
     $\cdot$ epochs 5\\
     10\,000 steps $\times$ seeds \{7, 19, 42\}};

  \draw[arr] (inject.south) -- ++(0,-0.4) -| (ppo.north);

  % Row 3
  \node[box=9cm, fill=green!10, below=1.2cm of ppo] (eval)
    {\textcircled{5} Evaluate on test window\\
     Replay policy $\to$ equity curve $\to$ Sharpe, MDD, preservation};

  \draw[arr] (ppo.south) -- (eval.north);

  % Note
  \node[below=0.5cm of eval, font=\footnotesize\itshape]
    {Baseline skips steps \textcircled{1}--\textcircled{3}:
     same PPO, no uncertainty channel.};
\end{tikzpicture}
\caption{End-to-end training pipeline from raw prices to evaluation metrics.}
\label{fig:pipeline}
\end{figure}


\section{WP1 --- Data and forecaster}
\label{sec:wp1}

\paragraph{T1.1 --- Data pipeline.} The function \texttt{fetch\_close\_frame} in \texttt{experiments/common.py} downloads daily adjusted-close prices via \texttt{yfinance} for any ticker and date range. The function is cached on disk so that subsequent calls do not re-hit the network. The protocol file \texttt{experiments/configs/dissertation\_protocol.json} contains the named ticker groups (\texttt{market\_sample}, \texttt{market\_sample\_stocks}, \texttt{market\_sample\_etfs}) and the train, validation and test split dates.

\paragraph{T1.2 --- Probabilistic LSTM.} The class \texttt{ProbabilisticLSTM} in \texttt{experiments/runners/run\_probabilistic\_agent.py} implements the architecture of Equation~\ref{eq:lstm_heads}: a two-layer LSTM with hidden dimension 32 and two linear heads. An optional dropout layer (probability $p_{\mathrm{drop}} = 0.1$) sits between the LSTM and the heads to support the epistemic mode of Section~\ref{sec:epistemic_mode}; in aleatoric mode the dropout probability is zero.

\paragraph{T1.3 --- Forecaster training.} Training uses the Gaussian negative log-likelihood loss of Equation~\ref{eq:gaussian_nll} and the Adam optimiser at learning rate $10^{-3}$ for 20 epochs. The training is fast --- under a minute on a single CPU core for any one ticker --- which is why the same script runs locally without a GPU.

\paragraph{T1.4 --- Uncertainty estimation.} The function \texttt{estimate\_uncertainty} produces the per-day uncertainty score $u_t$. The signature has an explicit \texttt{mode} argument that selects between aleatoric (Equation~\ref{eq:uncertainty_score}) and epistemic (Equation~\ref{eq:epistemic_uncertainty}), and a \texttt{mc\_passes} argument that controls the number of Monte-Carlo dropout passes in the epistemic mode (default 20). The returned array is min--max normalised to the unit interval.

\section{WP2 --- Environment and policy}
\label{sec:wp2}

\paragraph{T2.1 --- The \texttt{StockEnv} class.} The trading environment is implemented as the class \texttt{StockEnv} in \texttt{experiments/common.py}, exposing the standard Gymnasium API: \texttt{reset} returns the initial observation, and \texttt{step(action)} returns the next observation, the reward, a done flag, a truncated flag, and an info dictionary. The observation space is a flat float32 vector of dimension $L + 2$ (the lookback window of log-returns, the normalised position size and the uncertainty score, see Equation~\ref{eq:state}). The action space is a single continuous variable in $[-1, 1]$.

The trade-size computation inside \texttt{step} is the literal implementation of Equation~\ref{eq:probabilistic_trade_size}. When the environment is constructed without an uncertainty array (or with an all-zeros uncertainty array) the trade-scaling and risk-on-guard factors collapse to one and the environment behaves as the baseline-PPO environment of Equation~\ref{eq:baseline_trade_size}.

\paragraph{T2.2 --- Baseline PPO training.} The runner \texttt{experiments/runners/run\_baseline.py} loads the protocol, constructs a \texttt{StockEnv} with an all-zeros uncertainty array, instantiates a Stable-Baselines3 PPO model with an MLP policy and the hyper-parameters of Section~\ref{sec:baseline_ppo}, and trains for the requested number of time-steps. After training, the policy is evaluated deterministically on the test window and the resulting equity curve and metrics are written to \texttt{experiments/results/}.

\paragraph{T2.3 --- Probabilistic PPO training (Arm B, aleatoric).} The runner \texttt{experiments/runners/run\_probabilistic\_agent.py} extends the baseline by first training the probabilistic LSTM (T1.2 + T1.3), then computing the uncertainty score in aleatoric mode (T1.4 with \texttt{mode=\textquotesingle aleatoric\textquotesingle}), and finally training a PPO policy on the augmented \texttt{StockEnv}.

\paragraph{T2.4 --- Probabilistic PPO training (Arm C, epistemic).} The same runner with \texttt{-{}-uncertainty-mode epistemic} switches the uncertainty computation to Equation~\ref{eq:epistemic_uncertainty}. Two additional command-line flags expose the dropout probability and the number of Monte-Carlo passes: \texttt{-{}-mc-dropout 0.1 -{}-mc-passes 20}. The PPO training itself is unchanged; only the uncertainty array fed to \texttt{StockEnv} differs.

\section{WP3 --- Experimental grid}
\label{sec:wp3}

\paragraph{T3.1 --- The protocol file.} \texttt{experiments/configs/dissertation\_protocol.json} pins every reproducibility-critical choice: the date splits, the random seeds, the PPO hyper-parameters, the ticker groups, the uncertainty mode defaults, and the per-experiment training budgets. Every runner script reads from this file rather than from hard-coded constants. Changing one number in the protocol changes every downstream result.

\paragraph{T3.2 --- Single-fold runners.} The three runners in WP2 each accept the flags \texttt{-{}-tickers}, \texttt{-{}-seeds}, \texttt{-{}-timesteps} and \texttt{-{}-tag}, so that an arbitrary subset of the grid can be requested from the command line. Output files are tagged with the suffix so that successive runs do not overwrite each other.

\paragraph{T3.3 --- Walk-forward runner.} The runner \texttt{experiments/runners/run\_walk\_forward.py} iterates over the four protocol folds (\texttt{wf\_2018\_2019}, \texttt{wf\_2020\_2021}, \texttt{wf\_2022\_2023}, \texttt{wf\_2024\_2025}) and trains and evaluates per (ticker, fold, seed) cell. The Phase-1 walk-forward grid uses 4 tickers $\times$ 4 folds $\times$ 3 seeds = 48 cells per agent, which is CPU-feasible in one working day. The Phase-2 grid lifts this to 70 tickers $\times$ 4 folds $\times$ 10 seeds and is GPU-only.

\paragraph{T3.4 --- The Colab notebook.} \texttt{notebooks/02\_Full\_Experiments.ipynb} is the Colab-ready orchestrator for the Phase-2 grid. Section~\ref{sec:notebooks_as_evidence} below describes it in more detail.

\section{WP4 --- Evaluation and reporting}
\label{sec:wp4}

\paragraph{T4.1 --- Aggregation.} The script \texttt{experiments/aggregate\_results.py} collects every per-cell JSON file matching a tag suffix and produces median, inter-quartile range and win-count summaries across seeds. The output is a single CSV that drops directly into the tables of Chapter~\ref{ch:results} and Appendix~\ref{app:full_results}.

\paragraph{T4.2 --- Plotting.} The script \texttt{reports/builders/plot\_phase2\_charts.py} regenerates every Chapter~\ref{ch:results} figure deterministically from the committed Phase-2 CSVs in \texttt{experiments/results/}. The forecaster calibration figure is built separately by \texttt{reports/builders/build\_forecaster\_calibration.py}.

\paragraph{T4.3 --- Document builders.} The canonical dissertation is the LaTeX source at \texttt{latex/main.tex}, compiled with \texttt{pdflatex} and \texttt{biber} (see \texttt{latex/build.sh}). An editable Word copy is produced by \texttt{latex/build\_docx.sh}, which never modifies the LaTeX sources. Both document paths consume the same underlying CSVs and per-cell JSON produced by T4.1.

\section{Reproducibility}
\label{sec:reproducibility_implementation}

Three measures make the implementation reproducible. The first is the protocol file: every hyper-parameter that matters is in one JSON file and every runner reads from it. The second is the seed control: each runner sets the global random seed for Python, NumPy, PyTorch and Stable-Baselines3 at the top of \texttt{main} before any random operation. The third is the public repository: every script, every notebook and every protocol file lives at \url{https://github.com/TheFinix13/dissertation-project}, and Appendix~\ref{app:reproducibility} lists the exact command-line invocations that reproduce each table and figure in this dissertation.

\section{How the agent trades --- a day-by-day walkthrough}
\label{sec:day_by_day_walkthrough}

The tables and equity curves in Chapter~\ref{ch:results} report \emph{final} outcomes. To make the mechanism concrete, this section walks through two real test-window days end to end: a day on which the risk-on guard fired and a buy was blocked, and a day on which the trade-scaling factor capped exposure during a calm bull market and the agent gave up part of the move. Both are days that the day-by-day notebook (\texttt{notebooks/01\_Project\_Walkthrough.ipynb}) reproduces; the prices and uncertainty scores below are taken from that run.

\subsection{A blocked-trade day: SPY, 14 January 2022}
\label{sec:walkthrough_spy_blocked}

By the close of 14 January 2022 the probabilistic-aleatoric agent on SPY held 1\,254 shares, accumulated over the agent's slow build through the first two weeks of the test window. On that day SPY fell from \$464.53 to \$456.49, a drop of \$8.04 per share and a paper loss of roughly \$10\,000 across the holdings. The forecaster's normalised uncertainty score reached $u_t = 0.81$, just above the 80th-percentile threshold $\tau = 0.80$ that had been calibrated on the validation window (2019--2021).

The policy's raw action was $a_t = +0.45$. In the baseline-PPO environment of Equation~\ref{eq:baseline_trade_size} that action would have committed about \$50\,000 of remaining cash to additional SPY shares at the close, locking in further downside as the sell-off continued through the week of 18--24 January. In the probabilistic environment of Equation~\ref{eq:probabilistic_trade_size} the trade-scaling factor was already deflated to $\max(1 - 0.81, 0.10) = 0.19$, and the binary risk-on guard then zeroed the trade entirely because $a_t > 0$ and $u_t \geq \tau$. The agent did not buy. The position stayed at 1\,254 shares, the day's loss was the marked-to-market \$10\,000 (already incurred), and the additional \$50\,000 of capital was protected for the deeper drawdown that arrived in early February.

Two operational details are visible from this single day. First, the agent did not have to learn caution; the environment enforced it through the two extra factors in Equation~\ref{eq:probabilistic_trade_size}. Second, the guard is asymmetric --- if the action had been negative (a sell) the guard would have evaluated to one and the trade would have executed, so the agent retained the freedom to de-risk on a high-uncertainty day even while it lost the freedom to re-risk. This asymmetry is the architectural reason the policy can both ride a recovery and step out of a crash on the same threshold.

\subsection{A capped-exposure day: NVDA, 25 May 2023}
\label{sec:walkthrough_nvda_capped}

The mirror-image regime is a low-uncertainty bull move on a single-name trend stock. NVDA reported Q1 FY24 earnings on 24 May 2023 and gapped from a closing price near \$305 to an opening price near \$381 on 25 May 2023, a roughly $+24\,\%$ single-day move driven by record data-centre demand for AI accelerators. The forecaster, trained on the preceding decade in which moves of that magnitude were rare, emitted a low predictive variance on the days leading up to earnings; the normalised uncertainty score sat at $u_t \approx 0.18$ on the morning of 25 May, well below the threshold $\tau = 0.80$.

The risk-on guard therefore evaluated to one and did not interfere; the trade-scaling factor was $\max(1 - 0.18, 0.10) = 0.82$. The policy held a partial position --- about 82\,\% of the long exposure that the un-scaled trade size of Equation~\ref{eq:baseline_trade_size} would have committed --- because the trade-scaling factor caps exposure in proportion to forecast uncertainty even when the binary guard is silent. On a calm forecast day this is a feature: it keeps the agent from over-committing. On the morning of a $+24\,\%$ rip it is a cost. The agent participated in the move but at 82\,\% of buy-and-hold's notional, so the gap between the probabilistic agent's NVDA terminal value (\$7\,345\,000) and buy-and-hold's (\$8\,290\,000) was widened by another large step on this single day. Section~\ref{sec:wins_and_losses} returns to this regime as one of the two named failure modes of the design.

\subsection{What the walkthrough adds to the aggregate tables}

The blocked-trade day and the capped-exposure day are mirror illustrations of the trade-off the dissertation actually studies. On a chaotic day the binary guard suppresses a costly buy and protects capital; on a calm-trending day the same architecture costs the right tail of the realised return distribution. Neither pattern is visible from a summary statistic --- they only become legible when the reader watches the agent decide on one day with the actual prices, the actual uncertainty score, and the actual decision. The day-by-day notebook reproduces both walkthroughs in full and prints one row per day for the first sixty trading days of the test window.

\section{The notebooks as runnable evidence}
\label{sec:notebooks_as_evidence}

Two Jupyter notebooks live in \texttt{notebooks/} and are the dissertation's runnable evidence.

\paragraph{\texttt{01\_Project\_Walkthrough.ipynb}} runs locally on a single CPU in about five minutes. It walks through the whole pipeline end to end on SPY: it fetches the data, trains the probabilistic LSTM in both aleatoric and epistemic modes, trains all three reinforcement-learning arms (Baseline, Aleatoric, Epistemic) at the Phase-1 budget, reports the headline comparison table, plots the equity curves, and finishes with the day-by-day walkthrough of Section~\ref{sec:day_by_day_walkthrough}.

\paragraph{\texttt{02\_Full\_Experiments.ipynb}} runs on Google Colab with a T4 or A100 GPU runtime. It first runs a five-minute smoke test on SPY to verify the pipeline, then runs the Phase-2 grid (70 tickers $\times$ 10 seeds $\times$ 3 arms $\times$ 50\,000 PPO time-steps) followed by the walk-forward grid. Total wall time on the T4 is approximately four to six hours. The final cell zips every JSON and CSV in \texttt{experiments/results/} for download.

The notebooks are the reproducibility apparatus: any examiner with a Google account and a browser can verify the Phase-1 numbers in five minutes and the Phase-2 numbers in a few hours, with no local installation required.
